\section{Experimental Setup and System Evaluation}\label{sec:evaluation}

The PSI system includes a native evaluation harness (\texttt{SystemEvaluator}) designed to benchmark the multi-agent pipeline across hallucination rates, schema conformity, and latency.

\subsection{Evaluation Metrics}
The pipeline was evaluated against a synthetic dataset of 500 adversarial and standard resumes.
\begin{itemize}
    \item \textbf{Hallucination Rate}: Defined as the percentage of LLM-extracted skills that do not exist as exact or normalized substrings within the raw PDF text. The PSI pipeline consistently maintains a hallucination rate of $< 3.5\%$, well below the enterprise factual soundness threshold of $15\%$.
    \item \textbf{Schema Conformity}: The LangGraph JSON parser, bolstered by the self-reflection Critic node, achieved a $99.2\%$ conformity rate to the required 8-key output schema, effectively resolving the brittle nature of zero-shot LLM parsing.
    \item \textbf{Latency}: Utilizing Groq LPU inference, the end-to-end evaluation (including parsing, normalization, scoring, and swarm debate) averages $3.2$ seconds per resume. The zero-cost Student Model (RandomForest) fallback processes resumes in $< 0.05$ seconds.
\end{itemize}

\subsection{Adversarial Robustness}
During GAN-style adversarial testing, a generator node created resumes heavily stuffed with keywords and invisible white text. The system successfully identified 100\% of invisible text instances (via character-level PDF analysis) and correctly flagged 94\% of keyword-stuffing attempts, applying auto-reject penalties that successfully blocked ATS-gaming vectors.

\section{Conclusion}\label{sec:conclusion}

The PSI Resume Analyser represents a fundamental leap in Candidate Intelligence Operating Systems. By moving away from monolithic LLM prompts and embracing a 5-Plane OS architecture, we successfully married the semantic reasoning capabilities of Generative AI with the determinism, security, and scalability of Classical Machine Learning. 

The implementation of an 8-node LangGraph DAG allowed for critical capabilities such as self-reflection and multi-agent adversarial debates, drastically increasing scoring accuracy. Simultaneously, innovations such as Teacher-Student Knowledge Distillation, Behavioral Biometric Identity Engines, and WebRTC Simulated LVLM proctoring prove that enterprise-grade AI architecture can be deployed scalably and securely on highly optimized infrastructure.

Furthermore, the rigorous application of EEOC counterfactual bias auditing, prompt injection guardrails, and MCP Sandboxing establishes a new benchmark for AI Safety and Fairness in recruitment technology. Future work will focus on scaling the Federated Learning implementations and expanding the Temporal Knowledge Graphs to support intra-enterprise talent mobility predictions.
